
\documentclass{article}
\usepackage{colortbl}
\usepackage{makecell}
\usepackage{multirow}
\usepackage{supertabular}

\begin{document}

\newcounter{utterance}

\centering \large Interaction Transcript for game `dond', experiment `coop\_en', episode 13 with Qwen3{-}32B{-}t0.0.
\vspace{24pt}

{ \footnotesize  \setcounter{utterance}{1}
\setlength{\tabcolsep}{0pt}
\begin{supertabular}{c@{$\;$}|p{.15\linewidth}@{}p{.15\linewidth}p{.15\linewidth}p{.15\linewidth}p{.15\linewidth}p{.15\linewidth}}
    \# & $\;$A & \multicolumn{4}{c}{Game Master} & $\;\:$B\\
    \hline

    \theutterance \stepcounter{utterance}  
    & & \multicolumn{4}{p{0.6\linewidth}}{
        \cellcolor[rgb]{0.9,0.9,0.9}{
            \makecell[{{p{\linewidth}}}]{
                \texttt{\tiny{[P1$\langle$GM]}}
                \texttt{You are playing a cooperative negotiation game in which you have to agree on how to divide a set of items among you and another player.} \\
\\ 
\texttt{Rules:} \\
\texttt{(a) You and the other player are given a set of items. Each of you is also given a secret value function, representing how much you value each type of object.} \\
\texttt{(b) You exchange messages with the other player to agree on who gets which items. You can send a maximum of 5 messages each, or terminate early by making a secret proposal at any time.} \\
\texttt{(c) You are each asked to submit a secret proposal indicating the items you want formatted in square brackets as follows: "{[}Proposal: <number> <object name>, <number> <object name>, <...>{]}"} \\
\texttt{(d) If your proposals are complementary, i.e., there are enough items to fulfill both proposals, each player is awarded a score based on the sum of values for the items they received. Otherwise, both of you get zero points.} \\
\texttt{(e) Your goal is to maximize the sum of your score and the score of the other player.} \\
\\ 
\texttt{Let us start.} \\
\\ 
\texttt{The set of available items is:} \\
\\ 
\texttt{1 pen, 1 flower, 3 frames} \\
\\ 
\texttt{Your secret value function is:} \\
\\ 
\texttt{pen: 4, flower: 0, frame: 2} \\
\\ 
\texttt{IMPORTANT: Your messages, unless it is the secret proposal, are directly transmitted to the other player, so do not include any response to the rules or text announcing your message. To make a secret proposal, use the indicated format. Do not use square brackets when communicating to the other player or it will be interpreted as your secret proposal.} \\
            }
        }
    }
    & \\ \\

    \theutterance \stepcounter{utterance}  
    & \multicolumn{4}{p{0.6\linewidth}}{
        \cellcolor[rgb]{0.9,0.9,0.9}{
            \makecell[{{p{\linewidth}}}]{
                \texttt{\tiny{[P1$\rangle$GM]}}
                \texttt{<think>} \\
\texttt{Okay, let's see. I need to divide the items between me and another player. The items are 1 pen, 1 flower, and 3 frames. My value function is pen: 4, flower: 0, frame: 2. The other player's values are secret, so I don't know what they value. My goal is to maximize the total score (mine plus theirs) by making sure our proposals are complementary.} \\
\\ 
\texttt{First, I should figure out what the other player might value. Since I don't have their value function, I need to make some assumptions or ask for information. But since I can only send 5 messages, I need to be efficient. Maybe I can ask them what they value most. But wait, the messages are directly transmitted, so I can't include any meta instructions. So I should just ask them to communicate their preferences.} \\
\\ 
\texttt{But maybe I can think about possible scenarios. Let's consider possible value functions for the other player. For example, if they value the pen highly, they might want it, but since I value it at 4, which is high for me, I might want to get it. Alternatively, if they value the flower, but I value it at 0, so maybe they want it. Frames are 3 in total, each valued at 2 by me. So if the other player also values frames, we need to split them.} \\
\\ 
\texttt{But since the total score is the sum of both, I need to maximize the total. Let's think about possible allocations. For example, if I take the pen and some frames, and the other player takes the flower and the remaining frames. But since I don't value the flower, maybe they do. Alternatively, if they want the pen, maybe I should give it to them and take the frames. But I need to know their preferences.} \\
\\ 
\texttt{So the first step is to communicate with the other player to find out their preferences. Let me send a message asking them what they value most. For example, "What do you value most: pen, flower, or frames?" That way, they can respond and I can adjust my proposal accordingly.} \\
\\ 
\texttt{Alternatively, maybe I can propose a division based on possible scenarios. For example, if they value the pen, I can offer them the pen and some frames, and take the rest. But without knowing their values, it's risky. If I make a proposal without knowing, it might not be complementary} \\
            }
        }
    }
    & & \\ \\

    \theutterance \stepcounter{utterance}  
    & & & \multicolumn{2}{p{0.3\linewidth}}{
        \cellcolor[rgb]{0.9,0.9,0.9}{
            \makecell[{{p{\linewidth}}}]{
                \texttt{\tiny{[GM$|$GM]}}
                \texttt{continue} \\
            }
        }
    }
    & & \\ \\

    \theutterance \stepcounter{utterance}  
    & & & \multicolumn{2}{p{0.3\linewidth}}{
        \cellcolor[rgb]{0.9,0.9,0.9}{
            \makecell[{{p{\linewidth}}}]{
                \texttt{\tiny{[GM$|$GM]}}
                \texttt{<think>} \\
\texttt{Okay, let's see. I need to divide the items between me and another player. The items are 1 pen, 1 flower, and 3 frames. My value function is pen: 4, flower: 0, frame: 2. The other player's values are secret, so I don't know what they value. My goal is to maximize the total score (mine plus theirs) by making sure our proposals are complementary.} \\
\\ 
\texttt{First, I should figure out what the other player might value. Since I don't have their value function, I need to make some assumptions or ask for information. But since I can only send 5 messages, I need to be efficient. Maybe I can ask them what they value most. But wait, the messages are directly transmitted, so I can't include any meta instructions. So I should just ask them to communicate their preferences.} \\
\\ 
\texttt{But maybe I can think about possible scenarios. Let's consider possible value functions for the other player. For example, if they value the pen highly, they might want it, but since I value it at 4, which is high for me, I might want to get it. Alternatively, if they value the flower, but I value it at 0, so maybe they want it. Frames are 3 in total, each valued at 2 by me. So if the other player also values frames, we need to split them.} \\
\\ 
\texttt{But since the total score is the sum of both, I need to maximize the total. Let's think about possible allocations. For example, if I take the pen and some frames, and the other player takes the flower and the remaining frames. But since I don't value the flower, maybe they do. Alternatively, if they want the pen, maybe I should give it to them and take the frames. But I need to know their preferences.} \\
\\ 
\texttt{So the first step is to communicate with the other player to find out their preferences. Let me send a message asking them what they value most. For example, "What do you value most: pen, flower, or frames?" That way, they can respond and I can adjust my proposal accordingly.} \\
\\ 
\texttt{Alternatively, maybe I can propose a division based on possible scenarios. For example, if they value the pen, I can offer them the pen and some frames, and take the rest. But without knowing their values, it's risky. If I make a proposal without knowing, it might not be complementary} \\
            }
        }
    }
    & & \\ \\

    \theutterance \stepcounter{utterance}  
    & & \multicolumn{4}{p{0.6\linewidth}}{
        \cellcolor[rgb]{0.9,0.9,0.9}{
            \makecell[{{p{\linewidth}}}]{
                \texttt{\tiny{[GM$\rangle$P2]}}
                \texttt{You are playing a cooperative negotiation game in which you have to agree on how to divide a set of items among you and another player.} \\
\\ 
\texttt{Rules:} \\
\texttt{(a) You and the other player are given a set of items. Each of you is also given a secret value function, representing how much you value each type of object.} \\
\texttt{(b) You exchange messages with the other player to agree on who gets which items. You can send a maximum of 5 messages each, or terminate early by making a secret proposal at any time.} \\
\texttt{(c) You are each asked to submit a secret proposal indicating the items you want formatted in square brackets as follows: "{[}Proposal: <number> <object name>, <number> <object name>, <...>{]}"} \\
\texttt{(d) If your proposals are complementary, i.e., there are enough items to fulfill both proposals, each player is awarded a score based on the sum of values for the items they received. Otherwise, both of you get zero points.} \\
\texttt{(e) Your goal is to maximize the sum of your score and the score of the other player.} \\
\\ 
\texttt{Let us start.} \\
\\ 
\texttt{The set of available items is:} \\
\\ 
\texttt{1 pen, 1 flower, 3 frames} \\
\\ 
\texttt{Your secret value function is:} \\
\\ 
\texttt{pen: 2, flower: 5, frame: 1} \\
\\ 
\texttt{IMPORTANT: Your messages, unless it is the secret proposal, are directly transmitted to the other player, so do not include any response to the rules or text announcing your message. To make a secret proposal, use the indicated format. Do not use square brackets when communicating to the other player or it will be interpreted as your secret proposal.} \\
\\ 
\\ 
\texttt{<think>} \\
\texttt{Okay, let's see. I need to divide the items between me and another player. The items are 1 pen, 1 flower, and 3 frames. My value function is pen: 4, flower: 0, frame: 2. The other player's values are secret, so I don't know what they value. My goal is to maximize the total score (mine plus theirs) by making sure our proposals are complementary.} \\
\\ 
\texttt{First, I should figure out what the other player might value. Since I don't have their value function, I need to make some assumptions or ask for information. But since I can only send 5 messages, I need to be efficient. Maybe I can ask them what they value most. But wait, the messages are directly transmitted, so I can't include any meta instructions. So I should just ask them to communicate their preferences.} \\
\\ 
\texttt{But maybe I can think about possible scenarios. Let's consider possible value functions for the other player. For example, if they value the pen highly, they might want it, but since I value it at 4, which is high for me, I might want to get it. Alternatively, if they value the flower, but I value it at 0, so maybe they want it. Frames are 3 in total, each valued at 2 by me. So if the other player also values frames, we need to split them.} \\
\\ 
\texttt{But since the total score is the sum of both, I need to maximize the total. Let's think about possible allocations. For example, if I take the pen and some frames, and the other player takes the flower and the remaining frames. But since I don't value the flower, maybe they do. Alternatively, if they want the pen, maybe I should give it to them and take the frames. But I need to know their preferences.} \\
\\ 
\texttt{So the first step is to communicate with the other player to find out their preferences. Let me send a message asking them what they value most. For example, "What do you value most: pen, flower, or frames?" That way, they can respond and I can adjust my proposal accordingly.} \\
\\ 
\texttt{Alternatively, maybe I can propose a division based on possible scenarios. For example, if they value the pen, I can offer them the pen and some frames, and take the rest. But without knowing their values, it's risky. If I make a proposal without knowing, it might not be complementary} \\
            }
        }
    }
    & \\ \\

    \theutterance \stepcounter{utterance}  
    & & & \multicolumn{4}{p{0.6\linewidth}}{
        \cellcolor[rgb]{0.9,0.9,0.9}{
            \makecell[{{p{\linewidth}}}]{
                \texttt{\tiny{[GM$\langle$P2]}}
                \texttt{<think>} \\
\texttt{Okay, let's see. The user is playing a cooperative negotiation game where they need to divide items between themselves and another player. The items are 1 pen, 1 flower, and 3 frames. Their own value function is pen: 4, flower: 0, frame: 2. The other player's values are secret. The goal is to maximize the total score (the sum of both players' scores) by making complementary proposals.} \\
\\ 
\texttt{First, the user needs to figure out the other player's preferences. Since they can send up to 5 messages, they should use that to gather information. The key is to ask the other player about their valuations. For example, asking which items they value most. But since messages are direct and can't include meta instructions, the user should phrase it as a question like, "What do you value most: pen, flower, or frames?" } \\
\\ 
\texttt{Once they get a response, they can adjust their proposal. Suppose the other player says they value the pen highly. Then the user might offer the pen to the other player and take the frames and flower. But since the user values the pen at 4, which is high, maybe they should keep it if the other player doesn't value it. Alternatively, if the other player values the flower, the user can take the pen and frames, leaving the flower for the other. } \\
\\ 
\texttt{Frames are 3 in total. If both players value frames, they need to split them. The user's value per frame is 2, so 3 frames would give 6. If the other player also values frames, maybe they can split 2 and 1. But without knowing their value, it's tricky. } \\
\\ 
\texttt{The user should also consider possible scenarios. For example, if the other player's values are pen: high, flower: low, frames: medium, then the optimal allocation would be to give the pen to the other, the flower to the user (but user values it 0, so maybe not), and split frames. Wait, the user values the flower at 0, so maybe the other player wants it. } \\
\\ 
\texttt{The main challenge is uncertainty about the other player's values. The user needs to communicate to get that info. If they can't, they might have to make an educated guess. For instance, if the other player's values are not known, the safest bet is to propose taking the pen (since they value it 4} \\
            }
        }
    }
     \\ \\

    \theutterance \stepcounter{utterance}  
    & & & \multicolumn{2}{p{0.3\linewidth}}{
        \cellcolor[rgb]{0.9,0.9,0.9}{
            \makecell[{{p{\linewidth}}}]{
                \texttt{\tiny{[GM$|$GM]}}
                \texttt{continue} \\
            }
        }
    }
    & & \\ \\

    \theutterance \stepcounter{utterance}  
    & & & \multicolumn{2}{p{0.3\linewidth}}{
        \cellcolor[rgb]{0.9,0.9,0.9}{
            \makecell[{{p{\linewidth}}}]{
                \texttt{\tiny{[GM$|$GM]}}
                \texttt{<think>} \\
\texttt{Okay, let's see. The user is playing a cooperative negotiation game where they need to divide items between themselves and another player. The items are 1 pen, 1 flower, and 3 frames. Their own value function is pen: 4, flower: 0, frame: 2. The other player's values are secret. The goal is to maximize the total score (the sum of both players' scores) by making complementary proposals.} \\
\\ 
\texttt{First, the user needs to figure out the other player's preferences. Since they can send up to 5 messages, they should use that to gather information. The key is to ask the other player about their valuations. For example, asking which items they value most. But since messages are direct and can't include meta instructions, the user should phrase it as a question like, "What do you value most: pen, flower, or frames?" } \\
\\ 
\texttt{Once they get a response, they can adjust their proposal. Suppose the other player says they value the pen highly. Then the user might offer the pen to the other player and take the frames and flower. But since the user values the pen at 4, which is high, maybe they should keep it if the other player doesn't value it. Alternatively, if the other player values the flower, the user can take the pen and frames, leaving the flower for the other. } \\
\\ 
\texttt{Frames are 3 in total. If both players value frames, they need to split them. The user's value per frame is 2, so 3 frames would give 6. If the other player also values frames, maybe they can split 2 and 1. But without knowing their value, it's tricky. } \\
\\ 
\texttt{The user should also consider possible scenarios. For example, if the other player's values are pen: high, flower: low, frames: medium, then the optimal allocation would be to give the pen to the other, the flower to the user (but user values it 0, so maybe not), and split frames. Wait, the user values the flower at 0, so maybe the other player wants it. } \\
\\ 
\texttt{The main challenge is uncertainty about the other player's values. The user needs to communicate to get that info. If they can't, they might have to make an educated guess. For instance, if the other player's values are not known, the safest bet is to propose taking the pen (since they value it 4} \\
            }
        }
    }
    & & \\ \\

    \theutterance \stepcounter{utterance}  
    & & \multicolumn{4}{p{0.6\linewidth}}{
        \cellcolor[rgb]{0.9,0.9,0.9}{
            \makecell[{{p{\linewidth}}}]{
                \texttt{\tiny{[P1$\langle$GM]}}
                \texttt{<think>} \\
\texttt{Okay, let's see. The user is playing a cooperative negotiation game where they need to divide items between themselves and another player. The items are 1 pen, 1 flower, and 3 frames. Their own value function is pen: 4, flower: 0, frame: 2. The other player's values are secret. The goal is to maximize the total score (the sum of both players' scores) by making complementary proposals.} \\
\\ 
\texttt{First, the user needs to figure out the other player's preferences. Since they can send up to 5 messages, they should use that to gather information. The key is to ask the other player about their valuations. For example, asking which items they value most. But since messages are direct and can't include meta instructions, the user should phrase it as a question like, "What do you value most: pen, flower, or frames?" } \\
\\ 
\texttt{Once they get a response, they can adjust their proposal. Suppose the other player says they value the pen highly. Then the user might offer the pen to the other player and take the frames and flower. But since the user values the pen at 4, which is high, maybe they should keep it if the other player doesn't value it. Alternatively, if the other player values the flower, the user can take the pen and frames, leaving the flower for the other. } \\
\\ 
\texttt{Frames are 3 in total. If both players value frames, they need to split them. The user's value per frame is 2, so 3 frames would give 6. If the other player also values frames, maybe they can split 2 and 1. But without knowing their value, it's tricky. } \\
\\ 
\texttt{The user should also consider possible scenarios. For example, if the other player's values are pen: high, flower: low, frames: medium, then the optimal allocation would be to give the pen to the other, the flower to the user (but user values it 0, so maybe not), and split frames. Wait, the user values the flower at 0, so maybe the other player wants it. } \\
\\ 
\texttt{The main challenge is uncertainty about the other player's values. The user needs to communicate to get that info. If they can't, they might have to make an educated guess. For instance, if the other player's values are not known, the safest bet is to propose taking the pen (since they value it 4} \\
            }
        }
    }
    & \\ \\

    \theutterance \stepcounter{utterance}  
    & \multicolumn{4}{p{0.6\linewidth}}{
        \cellcolor[rgb]{0.9,0.9,0.9}{
            \makecell[{{p{\linewidth}}}]{
                \texttt{\tiny{[P1$\rangle$GM]}}
                \texttt{<think>} \\
\texttt{Okay, let's see. I need to figure out how to divide the items between me and the other player. The items are 1 pen, 1 flower, and 3 frames. My values are pen:4, flower:0, frame:2. The other player's values are secret, so I don't know what they want. My goal is to maximize the total score (mine plus theirs) by making sure our proposals are complementary.} \\
\\ 
\texttt{First, I should probably ask the other player what they value most. Since I can send up to 5 messages, I can use that to get some information. Let me start by asking them which items they value the most. For example, "What do you value most: pen, flower, or frames?" That way, they can respond and I can adjust my proposal accordingly.} \\
\\ 
\texttt{If they say they value the pen, then I might want to give them the pen and take the frames and flower. But wait, I value the pen at 4, which is pretty high. If they also value the pen, maybe we need to negotiate. But if they don't value the pen, maybe I should take it. Alternatively, if they value the flower, which I don't care about, I can take the pen and frames, and let them have the flower.} \\
\\ 
\texttt{Frames are 3 in total. If both of us value frames, we need to split them. Since I get 2 per frame, 3 frames would give me 6. If the other player also values frames, maybe we can split them 2 and 1. But without knowing their value, it's hard to say. Maybe I should propose taking 2 frames and the pen, leaving 1 frame and the flower for them. But if they value the flower, that might work. If they don't, maybe they want the frames.} \\
\\ 
\texttt{Wait, the flower is only 1, and I value it at 0. So if the other player values the flower, they might want it. So maybe I should take the pen and all 3 frames, and let them have the flower. That way, my score would be 4 (pen) + 3*2 (frames) = 10. If they value the flower, their score would be whatever they value it at. But if they don't value the flower, maybe they want the frames. Hmm.} \\
\\ 
\texttt{Alternatively, if they value the frames more than the} \\
            }
        }
    }
    & & \\ \\

    \theutterance \stepcounter{utterance}  
    & & & \multicolumn{2}{p{0.3\linewidth}}{
        \cellcolor[rgb]{0.9,0.9,0.9}{
            \makecell[{{p{\linewidth}}}]{
                \texttt{\tiny{[GM$|$GM]}}
                \texttt{continue} \\
            }
        }
    }
    & & \\ \\

    \theutterance \stepcounter{utterance}  
    & & & \multicolumn{2}{p{0.3\linewidth}}{
        \cellcolor[rgb]{0.9,0.9,0.9}{
            \makecell[{{p{\linewidth}}}]{
                \texttt{\tiny{[GM$|$GM]}}
                \texttt{<think>} \\
\texttt{Okay, let's see. I need to figure out how to divide the items between me and the other player. The items are 1 pen, 1 flower, and 3 frames. My values are pen:4, flower:0, frame:2. The other player's values are secret, so I don't know what they want. My goal is to maximize the total score (mine plus theirs) by making sure our proposals are complementary.} \\
\\ 
\texttt{First, I should probably ask the other player what they value most. Since I can send up to 5 messages, I can use that to get some information. Let me start by asking them which items they value the most. For example, "What do you value most: pen, flower, or frames?" That way, they can respond and I can adjust my proposal accordingly.} \\
\\ 
\texttt{If they say they value the pen, then I might want to give them the pen and take the frames and flower. But wait, I value the pen at 4, which is pretty high. If they also value the pen, maybe we need to negotiate. But if they don't value the pen, maybe I should take it. Alternatively, if they value the flower, which I don't care about, I can take the pen and frames, and let them have the flower.} \\
\\ 
\texttt{Frames are 3 in total. If both of us value frames, we need to split them. Since I get 2 per frame, 3 frames would give me 6. If the other player also values frames, maybe we can split them 2 and 1. But without knowing their value, it's hard to say. Maybe I should propose taking 2 frames and the pen, leaving 1 frame and the flower for them. But if they value the flower, that might work. If they don't, maybe they want the frames.} \\
\\ 
\texttt{Wait, the flower is only 1, and I value it at 0. So if the other player values the flower, they might want it. So maybe I should take the pen and all 3 frames, and let them have the flower. That way, my score would be 4 (pen) + 3*2 (frames) = 10. If they value the flower, their score would be whatever they value it at. But if they don't value the flower, maybe they want the frames. Hmm.} \\
\\ 
\texttt{Alternatively, if they value the frames more than the} \\
            }
        }
    }
    & & \\ \\

    \theutterance \stepcounter{utterance}  
    & & \multicolumn{4}{p{0.6\linewidth}}{
        \cellcolor[rgb]{0.9,0.9,0.9}{
            \makecell[{{p{\linewidth}}}]{
                \texttt{\tiny{[GM$\rangle$P2]}}
                \texttt{<think>} \\
\texttt{Okay, let's see. I need to figure out how to divide the items between me and the other player. The items are 1 pen, 1 flower, and 3 frames. My values are pen:4, flower:0, frame:2. The other player's values are secret, so I don't know what they want. My goal is to maximize the total score (mine plus theirs) by making sure our proposals are complementary.} \\
\\ 
\texttt{First, I should probably ask the other player what they value most. Since I can send up to 5 messages, I can use that to get some information. Let me start by asking them which items they value the most. For example, "What do you value most: pen, flower, or frames?" That way, they can respond and I can adjust my proposal accordingly.} \\
\\ 
\texttt{If they say they value the pen, then I might want to give them the pen and take the frames and flower. But wait, I value the pen at 4, which is pretty high. If they also value the pen, maybe we need to negotiate. But if they don't value the pen, maybe I should take it. Alternatively, if they value the flower, which I don't care about, I can take the pen and frames, and let them have the flower.} \\
\\ 
\texttt{Frames are 3 in total. If both of us value frames, we need to split them. Since I get 2 per frame, 3 frames would give me 6. If the other player also values frames, maybe we can split them 2 and 1. But without knowing their value, it's hard to say. Maybe I should propose taking 2 frames and the pen, leaving 1 frame and the flower for them. But if they value the flower, that might work. If they don't, maybe they want the frames.} \\
\\ 
\texttt{Wait, the flower is only 1, and I value it at 0. So if the other player values the flower, they might want it. So maybe I should take the pen and all 3 frames, and let them have the flower. That way, my score would be 4 (pen) + 3*2 (frames) = 10. If they value the flower, their score would be whatever they value it at. But if they don't value the flower, maybe they want the frames. Hmm.} \\
\\ 
\texttt{Alternatively, if they value the frames more than the} \\
            }
        }
    }
    & \\ \\

    \theutterance \stepcounter{utterance}  
    & & & \multicolumn{4}{p{0.6\linewidth}}{
        \cellcolor[rgb]{0.9,0.9,0.9}{
            \makecell[{{p{\linewidth}}}]{
                \texttt{\tiny{[GM$\langle$P2]}}
                \texttt{<think>} \\
\texttt{Okay, let's see. I need to figure out how to divide the items between me and the other player. The items are 1 pen, 1 flower, and 3 frames. My values are pen:4, flower:0, frame:2. The other player's values are secret, so I don't know what they want. My goal is to maximize the total score (mine plus theirs) by making sure our proposals are complementary.} \\
\\ 
\texttt{First, I should probably ask the other player what they value most. Since I can send up to 5 messages, I can use that to get some information. Let me start by asking them which items they value the most. For example, "What do you value most: pen, flower, or frames?" That way, they can respond and I can adjust my proposal accordingly.} \\
\\ 
\texttt{If they say they value the pen, then I might want to give them the pen and take the frames and flower. But wait, I value the pen at 4, which is pretty high. If they also value the pen, maybe we need to negotiate. But if they don't value the pen, maybe I should take it. Alternatively, if they value the flower, which I don't care about, I can take the pen and frames, and let them have the flower.} \\
\\ 
\texttt{Frames are 3 in total. If both of us value frames, we need to split them. Since I get 2 per frame, 3 frames would give me 6. If the other player also values frames, maybe we can split them 2 and 1. But without knowing their value, it's hard to say. Maybe I should propose taking 2 frames and the pen, leaving 1 frame and the flower for them. But if they value the flower, that might work. If they don't, maybe they want the frames.} \\
\\ 
\texttt{Wait, the flower is only 1, and I value it at 0. So if the other player values the flower, they might want it. So maybe I should take the pen and all 3 frames, and let them have the flower. That way, my score would be 4 (pen) + 3*2 (frames) = 10. If they value the flower, their score would be whatever they value it at. But if they don't value the flower, maybe they want the frames. Hmm.} \\
\\ 
\texttt{Alternatively, if they value the frames more than the} \\
            }
        }
    }
     \\ \\

    \theutterance \stepcounter{utterance}  
    & & & \multicolumn{2}{p{0.3\linewidth}}{
        \cellcolor[rgb]{0.9,0.9,0.9}{
            \makecell[{{p{\linewidth}}}]{
                \texttt{\tiny{[GM$|$GM]}}
                \texttt{continue} \\
            }
        }
    }
    & & \\ \\

    \theutterance \stepcounter{utterance}  
    & & & \multicolumn{2}{p{0.3\linewidth}}{
        \cellcolor[rgb]{0.9,0.9,0.9}{
            \makecell[{{p{\linewidth}}}]{
                \texttt{\tiny{[GM$|$GM]}}
                \texttt{<think>} \\
\texttt{Okay, let's see. I need to figure out how to divide the items between me and the other player. The items are 1 pen, 1 flower, and 3 frames. My values are pen:4, flower:0, frame:2. The other player's values are secret, so I don't know what they want. My goal is to maximize the total score (mine plus theirs) by making sure our proposals are complementary.} \\
\\ 
\texttt{First, I should probably ask the other player what they value most. Since I can send up to 5 messages, I can use that to get some information. Let me start by asking them which items they value the most. For example, "What do you value most: pen, flower, or frames?" That way, they can respond and I can adjust my proposal accordingly.} \\
\\ 
\texttt{If they say they value the pen, then I might want to give them the pen and take the frames and flower. But wait, I value the pen at 4, which is pretty high. If they also value the pen, maybe we need to negotiate. But if they don't value the pen, maybe I should take it. Alternatively, if they value the flower, which I don't care about, I can take the pen and frames, and let them have the flower.} \\
\\ 
\texttt{Frames are 3 in total. If both of us value frames, we need to split them. Since I get 2 per frame, 3 frames would give me 6. If the other player also values frames, maybe we can split them 2 and 1. But without knowing their value, it's hard to say. Maybe I should propose taking 2 frames and the pen, leaving 1 frame and the flower for them. But if they value the flower, that might work. If they don't, maybe they want the frames.} \\
\\ 
\texttt{Wait, the flower is only 1, and I value it at 0. So if the other player values the flower, they might want it. So maybe I should take the pen and all 3 frames, and let them have the flower. That way, my score would be 4 (pen) + 3*2 (frames) = 10. If they value the flower, their score would be whatever they value it at. But if they don't value the flower, maybe they want the frames. Hmm.} \\
\\ 
\texttt{Alternatively, if they value the frames more than the} \\
            }
        }
    }
    & & \\ \\

    \theutterance \stepcounter{utterance}  
    & & \multicolumn{4}{p{0.6\linewidth}}{
        \cellcolor[rgb]{0.9,0.9,0.9}{
            \makecell[{{p{\linewidth}}}]{
                \texttt{\tiny{[P1$\langle$GM]}}
                \texttt{<think>} \\
\texttt{Okay, let's see. I need to figure out how to divide the items between me and the other player. The items are 1 pen, 1 flower, and 3 frames. My values are pen:4, flower:0, frame:2. The other player's values are secret, so I don't know what they want. My goal is to maximize the total score (mine plus theirs) by making sure our proposals are complementary.} \\
\\ 
\texttt{First, I should probably ask the other player what they value most. Since I can send up to 5 messages, I can use that to get some information. Let me start by asking them which items they value the most. For example, "What do you value most: pen, flower, or frames?" That way, they can respond and I can adjust my proposal accordingly.} \\
\\ 
\texttt{If they say they value the pen, then I might want to give them the pen and take the frames and flower. But wait, I value the pen at 4, which is pretty high. If they also value the pen, maybe we need to negotiate. But if they don't value the pen, maybe I should take it. Alternatively, if they value the flower, which I don't care about, I can take the pen and frames, and let them have the flower.} \\
\\ 
\texttt{Frames are 3 in total. If both of us value frames, we need to split them. Since I get 2 per frame, 3 frames would give me 6. If the other player also values frames, maybe we can split them 2 and 1. But without knowing their value, it's hard to say. Maybe I should propose taking 2 frames and the pen, leaving 1 frame and the flower for them. But if they value the flower, that might work. If they don't, maybe they want the frames.} \\
\\ 
\texttt{Wait, the flower is only 1, and I value it at 0. So if the other player values the flower, they might want it. So maybe I should take the pen and all 3 frames, and let them have the flower. That way, my score would be 4 (pen) + 3*2 (frames) = 10. If they value the flower, their score would be whatever they value it at. But if they don't value the flower, maybe they want the frames. Hmm.} \\
\\ 
\texttt{Alternatively, if they value the frames more than the} \\
            }
        }
    }
    & \\ \\

    \theutterance \stepcounter{utterance}  
    & \multicolumn{4}{p{0.6\linewidth}}{
        \cellcolor[rgb]{0.9,0.9,0.9}{
            \makecell[{{p{\linewidth}}}]{
                \texttt{\tiny{[P1$\rangle$GM]}}
                \texttt{<think>} \\
\texttt{Okay, let's see. I need to figure out how to divide the items between me and the other player. The items are 1 pen, 1 flower, and 3 frames. My values are pen:4, flower:0, frame:2. The other player's values are secret, so I don't know what they want. My goal is to maximize the total score (mine plus theirs) by making sure our proposals are complementary.} \\
\\ 
\texttt{First, I should probably ask the other player what they value most. Since I can send up to 5 messages, I can use that to get some information. Let me start by asking them which items they value the most. For example, "What do you value most: pen, flower, or frames?" That way, they can respond and I can adjust my proposal accordingly.} \\
\\ 
\texttt{If they say they value the pen, then I might want to give them the pen and take the frames and flower. But wait, I value the pen at 4, which is pretty high. If they also value the pen, maybe we need to negotiate. But if they don't value the pen, maybe I should take it. Alternatively, if they value the flower, which I don't care about, I can take the pen and frames, and let them have the flower.} \\
\\ 
\texttt{Frames are 3 in total. If both of us value frames, we need to split them. Since I get 2 per frame, 3 frames would give me 6. If the other player also values frames, maybe we can split them 2 and 1. But without knowing their value, it's hard to say. Maybe I should propose taking 2 frames and the pen, leaving 1 frame and the flower for them. But if they value the flower, that might work. If they don't, maybe they want the frames.} \\
\\ 
\texttt{Wait, the flower is only 1, and I value it at 0. So if the other player values the flower, they might want it. So maybe I should take the pen and all 3 frames, and let them have the flower. That way, my score would be 4 (pen) + 3*2 (frames) = 10. If they value the flower, their score would be whatever they value it at. But if they don't value the flower, maybe they want the frames. Hmm.} \\
\\ 
\texttt{Alternatively, if they value the frames more than the} \\
            }
        }
    }
    & & \\ \\

    \theutterance \stepcounter{utterance}  
    & & & \multicolumn{2}{p{0.3\linewidth}}{
        \cellcolor[rgb]{0.9,0.9,0.9}{
            \makecell[{{p{\linewidth}}}]{
                \texttt{\tiny{[GM$|$GM]}}
                \texttt{continue} \\
            }
        }
    }
    & & \\ \\

    \theutterance \stepcounter{utterance}  
    & & & \multicolumn{2}{p{0.3\linewidth}}{
        \cellcolor[rgb]{0.9,0.9,0.9}{
            \makecell[{{p{\linewidth}}}]{
                \texttt{\tiny{[GM$|$GM]}}
                \texttt{<think>} \\
\texttt{Okay, let's see. I need to figure out how to divide the items between me and the other player. The items are 1 pen, 1 flower, and 3 frames. My values are pen:4, flower:0, frame:2. The other player's values are secret, so I don't know what they want. My goal is to maximize the total score (mine plus theirs) by making sure our proposals are complementary.} \\
\\ 
\texttt{First, I should probably ask the other player what they value most. Since I can send up to 5 messages, I can use that to get some information. Let me start by asking them which items they value the most. For example, "What do you value most: pen, flower, or frames?" That way, they can respond and I can adjust my proposal accordingly.} \\
\\ 
\texttt{If they say they value the pen, then I might want to give them the pen and take the frames and flower. But wait, I value the pen at 4, which is pretty high. If they also value the pen, maybe we need to negotiate. But if they don't value the pen, maybe I should take it. Alternatively, if they value the flower, which I don't care about, I can take the pen and frames, and let them have the flower.} \\
\\ 
\texttt{Frames are 3 in total. If both of us value frames, we need to split them. Since I get 2 per frame, 3 frames would give me 6. If the other player also values frames, maybe we can split them 2 and 1. But without knowing their value, it's hard to say. Maybe I should propose taking 2 frames and the pen, leaving 1 frame and the flower for them. But if they value the flower, that might work. If they don't, maybe they want the frames.} \\
\\ 
\texttt{Wait, the flower is only 1, and I value it at 0. So if the other player values the flower, they might want it. So maybe I should take the pen and all 3 frames, and let them have the flower. That way, my score would be 4 (pen) + 3*2 (frames) = 10. If they value the flower, their score would be whatever they value it at. But if they don't value the flower, maybe they want the frames. Hmm.} \\
\\ 
\texttt{Alternatively, if they value the frames more than the} \\
            }
        }
    }
    & & \\ \\

    \theutterance \stepcounter{utterance}  
    & & \multicolumn{4}{p{0.6\linewidth}}{
        \cellcolor[rgb]{0.9,0.9,0.9}{
            \makecell[{{p{\linewidth}}}]{
                \texttt{\tiny{[GM$\rangle$P2]}}
                \texttt{<think>} \\
\texttt{Okay, let's see. I need to figure out how to divide the items between me and the other player. The items are 1 pen, 1 flower, and 3 frames. My values are pen:4, flower:0, frame:2. The other player's values are secret, so I don't know what they want. My goal is to maximize the total score (mine plus theirs) by making sure our proposals are complementary.} \\
\\ 
\texttt{First, I should probably ask the other player what they value most. Since I can send up to 5 messages, I can use that to get some information. Let me start by asking them which items they value the most. For example, "What do you value most: pen, flower, or frames?" That way, they can respond and I can adjust my proposal accordingly.} \\
\\ 
\texttt{If they say they value the pen, then I might want to give them the pen and take the frames and flower. But wait, I value the pen at 4, which is pretty high. If they also value the pen, maybe we need to negotiate. But if they don't value the pen, maybe I should take it. Alternatively, if they value the flower, which I don't care about, I can take the pen and frames, and let them have the flower.} \\
\\ 
\texttt{Frames are 3 in total. If both of us value frames, we need to split them. Since I get 2 per frame, 3 frames would give me 6. If the other player also values frames, maybe we can split them 2 and 1. But without knowing their value, it's hard to say. Maybe I should propose taking 2 frames and the pen, leaving 1 frame and the flower for them. But if they value the flower, that might work. If they don't, maybe they want the frames.} \\
\\ 
\texttt{Wait, the flower is only 1, and I value it at 0. So if the other player values the flower, they might want it. So maybe I should take the pen and all 3 frames, and let them have the flower. That way, my score would be 4 (pen) + 3*2 (frames) = 10. If they value the flower, their score would be whatever they value it at. But if they don't value the flower, maybe they want the frames. Hmm.} \\
\\ 
\texttt{Alternatively, if they value the frames more than the} \\
            }
        }
    }
    & \\ \\

    \theutterance \stepcounter{utterance}  
    & & & \multicolumn{4}{p{0.6\linewidth}}{
        \cellcolor[rgb]{0.9,0.9,0.9}{
            \makecell[{{p{\linewidth}}}]{
                \texttt{\tiny{[GM$\langle$P2]}}
                \texttt{<think>} \\
\texttt{Okay, let's see. I need to figure out how to divide the items between me and the other player. The items are 1 pen, 1 flower, and 3 frames. My values are pen:4, flower:0, frame:2. The other player's values are secret, so I don't know what they want. My goal is to maximize the total score (mine plus theirs) by making sure our proposals are complementary.} \\
\\ 
\texttt{First, I should probably ask the other player what they value most. Since I can send up to 5 messages, I can use that to get some information. Let me start by asking them which items they value the most. For example, "What do you value most: pen, flower, or frames?" That way, they can respond and I can adjust my proposal accordingly.} \\
\\ 
\texttt{If they say they value the pen, then I might want to give them the pen and take the frames and flower. But wait, I value the pen at 4, which is pretty high. If they also value the pen, maybe we need to negotiate. But if they don't value the pen, maybe I should take it. Alternatively, if they value the flower, which I don't care about, I can take the pen and frames, and let them have the flower.} \\
\\ 
\texttt{Frames are 3 in total. If both of us value frames, we need to split them. Since I get 2 per frame, 3 frames would give me 6. If the other player also values frames, maybe we can split them 2 and 1. But without knowing their value, it's hard to say. Maybe I should propose taking 2 frames and the pen, leaving 1 frame and the flower for them. But if they value the flower, that might work. If they don't, maybe they want the frames.} \\
\\ 
\texttt{Wait, the flower is only 1, and I value it at 0. So if the other player values the flower, they might want it. So maybe I should take the pen and all 3 frames, and let them have the flower. That way, my score would be 4 (pen) + 3*2 (frames) = 10. If they value the flower, their score would be whatever they value it at. But if they don't value the flower, maybe they want the frames. Hmm.} \\
\\ 
\texttt{Alternatively, if they value the frames more than the} \\
            }
        }
    }
     \\ \\

    \theutterance \stepcounter{utterance}  
    & & & \multicolumn{2}{p{0.3\linewidth}}{
        \cellcolor[rgb]{0.9,0.9,0.9}{
            \makecell[{{p{\linewidth}}}]{
                \texttt{\tiny{[GM$|$GM]}}
                \texttt{continue} \\
            }
        }
    }
    & & \\ \\

    \theutterance \stepcounter{utterance}  
    & & & \multicolumn{2}{p{0.3\linewidth}}{
        \cellcolor[rgb]{0.9,0.9,0.9}{
            \makecell[{{p{\linewidth}}}]{
                \texttt{\tiny{[GM$|$GM]}}
                \texttt{<think>} \\
\texttt{Okay, let's see. I need to figure out how to divide the items between me and the other player. The items are 1 pen, 1 flower, and 3 frames. My values are pen:4, flower:0, frame:2. The other player's values are secret, so I don't know what they want. My goal is to maximize the total score (mine plus theirs) by making sure our proposals are complementary.} \\
\\ 
\texttt{First, I should probably ask the other player what they value most. Since I can send up to 5 messages, I can use that to get some information. Let me start by asking them which items they value the most. For example, "What do you value most: pen, flower, or frames?" That way, they can respond and I can adjust my proposal accordingly.} \\
\\ 
\texttt{If they say they value the pen, then I might want to give them the pen and take the frames and flower. But wait, I value the pen at 4, which is pretty high. If they also value the pen, maybe we need to negotiate. But if they don't value the pen, maybe I should take it. Alternatively, if they value the flower, which I don't care about, I can take the pen and frames, and let them have the flower.} \\
\\ 
\texttt{Frames are 3 in total. If both of us value frames, we need to split them. Since I get 2 per frame, 3 frames would give me 6. If the other player also values frames, maybe we can split them 2 and 1. But without knowing their value, it's hard to say. Maybe I should propose taking 2 frames and the pen, leaving 1 frame and the flower for them. But if they value the flower, that might work. If they don't, maybe they want the frames.} \\
\\ 
\texttt{Wait, the flower is only 1, and I value it at 0. So if the other player values the flower, they might want it. So maybe I should take the pen and all 3 frames, and let them have the flower. That way, my score would be 4 (pen) + 3*2 (frames) = 10. If they value the flower, their score would be whatever they value it at. But if they don't value the flower, maybe they want the frames. Hmm.} \\
\\ 
\texttt{Alternatively, if they value the frames more than the} \\
            }
        }
    }
    & & \\ \\

    \theutterance \stepcounter{utterance}  
    & & \multicolumn{4}{p{0.6\linewidth}}{
        \cellcolor[rgb]{0.9,0.9,0.9}{
            \makecell[{{p{\linewidth}}}]{
                \texttt{\tiny{[P1$\langle$GM]}}
                \texttt{<think>} \\
\texttt{Okay, let's see. I need to figure out how to divide the items between me and the other player. The items are 1 pen, 1 flower, and 3 frames. My values are pen:4, flower:0, frame:2. The other player's values are secret, so I don't know what they want. My goal is to maximize the total score (mine plus theirs) by making sure our proposals are complementary.} \\
\\ 
\texttt{First, I should probably ask the other player what they value most. Since I can send up to 5 messages, I can use that to get some information. Let me start by asking them which items they value the most. For example, "What do you value most: pen, flower, or frames?" That way, they can respond and I can adjust my proposal accordingly.} \\
\\ 
\texttt{If they say they value the pen, then I might want to give them the pen and take the frames and flower. But wait, I value the pen at 4, which is pretty high. If they also value the pen, maybe we need to negotiate. But if they don't value the pen, maybe I should take it. Alternatively, if they value the flower, which I don't care about, I can take the pen and frames, and let them have the flower.} \\
\\ 
\texttt{Frames are 3 in total. If both of us value frames, we need to split them. Since I get 2 per frame, 3 frames would give me 6. If the other player also values frames, maybe we can split them 2 and 1. But without knowing their value, it's hard to say. Maybe I should propose taking 2 frames and the pen, leaving 1 frame and the flower for them. But if they value the flower, that might work. If they don't, maybe they want the frames.} \\
\\ 
\texttt{Wait, the flower is only 1, and I value it at 0. So if the other player values the flower, they might want it. So maybe I should take the pen and all 3 frames, and let them have the flower. That way, my score would be 4 (pen) + 3*2 (frames) = 10. If they value the flower, their score would be whatever they value it at. But if they don't value the flower, maybe they want the frames. Hmm.} \\
\\ 
\texttt{Alternatively, if they value the frames more than the} \\
            }
        }
    }
    & \\ \\

    \theutterance \stepcounter{utterance}  
    & \multicolumn{4}{p{0.6\linewidth}}{
        \cellcolor[rgb]{0.9,0.9,0.9}{
            \makecell[{{p{\linewidth}}}]{
                \texttt{\tiny{[P1$\rangle$GM]}}
                \texttt{<think>} \\
\texttt{Okay, let's see. I need to figure out how to divide the items between me and the other player. The items are 1 pen, 1 flower, and 3 frames. My values are pen:4, flower:0, frame:2. The other player's values are secret, so I don't know what they want. My goal is to maximize the total score (mine plus theirs) by making sure our proposals are complementary.} \\
\\ 
\texttt{First, I should probably ask the other player what they value most. Since I can send up to 5 messages, I can use that to get some information. Let me start by asking them which items they value the most. For example, "What do you value most: pen, flower, or frames?" That way, they can respond and I can adjust my proposal accordingly.} \\
\\ 
\texttt{If they say they value the pen, then I might want to give them the pen and take the frames and flower. But wait, I value the pen at 4, which is pretty high. If they also value the pen, maybe we need to negotiate. But if they don't value the pen, maybe I should take it. Alternatively, if they value the flower, which I don't care about, I can take the pen and frames, and let them have the flower.} \\
\\ 
\texttt{Frames are 3 in total. If both of us value frames, we need to split them. Since I get 2 per frame, 3 frames would give me 6. If the other player also values frames, maybe we can split them 2 and 1. But without knowing their value, it's hard to say. Maybe I should propose taking 2 frames and the pen, leaving 1 frame and the flower for them. But if they value the flower, that might work. If they don't, maybe they want the frames.} \\
\\ 
\texttt{Wait, the flower is only 1, and I value it at 0. So if the other player values the flower, they might want it. So maybe I should take the pen and all 3 frames, and let them have the flower. That way, my score would be 4 (pen) + 3*2 (frames) = 10. If they value the flower, their score would be whatever they value it at. But if they don't value the flower, maybe they want the frames. Hmm.} \\
\\ 
\texttt{Alternatively, if they value the frames more than the} \\
            }
        }
    }
    & & \\ \\

    \theutterance \stepcounter{utterance}  
    & & & \multicolumn{2}{p{0.3\linewidth}}{
        \cellcolor[rgb]{0.9,0.9,0.9}{
            \makecell[{{p{\linewidth}}}]{
                \texttt{\tiny{[GM$|$GM]}}
                \texttt{continue} \\
            }
        }
    }
    & & \\ \\

    \theutterance \stepcounter{utterance}  
    & & & \multicolumn{2}{p{0.3\linewidth}}{
        \cellcolor[rgb]{0.9,0.9,0.9}{
            \makecell[{{p{\linewidth}}}]{
                \texttt{\tiny{[GM$|$GM]}}
                \texttt{<think>} \\
\texttt{Okay, let's see. I need to figure out how to divide the items between me and the other player. The items are 1 pen, 1 flower, and 3 frames. My values are pen:4, flower:0, frame:2. The other player's values are secret, so I don't know what they want. My goal is to maximize the total score (mine plus theirs) by making sure our proposals are complementary.} \\
\\ 
\texttt{First, I should probably ask the other player what they value most. Since I can send up to 5 messages, I can use that to get some information. Let me start by asking them which items they value the most. For example, "What do you value most: pen, flower, or frames?" That way, they can respond and I can adjust my proposal accordingly.} \\
\\ 
\texttt{If they say they value the pen, then I might want to give them the pen and take the frames and flower. But wait, I value the pen at 4, which is pretty high. If they also value the pen, maybe we need to negotiate. But if they don't value the pen, maybe I should take it. Alternatively, if they value the flower, which I don't care about, I can take the pen and frames, and let them have the flower.} \\
\\ 
\texttt{Frames are 3 in total. If both of us value frames, we need to split them. Since I get 2 per frame, 3 frames would give me 6. If the other player also values frames, maybe we can split them 2 and 1. But without knowing their value, it's hard to say. Maybe I should propose taking 2 frames and the pen, leaving 1 frame and the flower for them. But if they value the flower, that might work. If they don't, maybe they want the frames.} \\
\\ 
\texttt{Wait, the flower is only 1, and I value it at 0. So if the other player values the flower, they might want it. So maybe I should take the pen and all 3 frames, and let them have the flower. That way, my score would be 4 (pen) + 3*2 (frames) = 10. If they value the flower, their score would be whatever they value it at. But if they don't value the flower, maybe they want the frames. Hmm.} \\
\\ 
\texttt{Alternatively, if they value the frames more than the} \\
            }
        }
    }
    & & \\ \\

    \theutterance \stepcounter{utterance}  
    & & \multicolumn{4}{p{0.6\linewidth}}{
        \cellcolor[rgb]{0.9,0.9,0.9}{
            \makecell[{{p{\linewidth}}}]{
                \texttt{\tiny{[GM$\rangle$P2]}}
                \texttt{<think>} \\
\texttt{Okay, let's see. I need to figure out how to divide the items between me and the other player. The items are 1 pen, 1 flower, and 3 frames. My values are pen:4, flower:0, frame:2. The other player's values are secret, so I don't know what they want. My goal is to maximize the total score (mine plus theirs) by making sure our proposals are complementary.} \\
\\ 
\texttt{First, I should probably ask the other player what they value most. Since I can send up to 5 messages, I can use that to get some information. Let me start by asking them which items they value the most. For example, "What do you value most: pen, flower, or frames?" That way, they can respond and I can adjust my proposal accordingly.} \\
\\ 
\texttt{If they say they value the pen, then I might want to give them the pen and take the frames and flower. But wait, I value the pen at 4, which is pretty high. If they also value the pen, maybe we need to negotiate. But if they don't value the pen, maybe I should take it. Alternatively, if they value the flower, which I don't care about, I can take the pen and frames, and let them have the flower.} \\
\\ 
\texttt{Frames are 3 in total. If both of us value frames, we need to split them. Since I get 2 per frame, 3 frames would give me 6. If the other player also values frames, maybe we can split them 2 and 1. But without knowing their value, it's hard to say. Maybe I should propose taking 2 frames and the pen, leaving 1 frame and the flower for them. But if they value the flower, that might work. If they don't, maybe they want the frames.} \\
\\ 
\texttt{Wait, the flower is only 1, and I value it at 0. So if the other player values the flower, they might want it. So maybe I should take the pen and all 3 frames, and let them have the flower. That way, my score would be 4 (pen) + 3*2 (frames) = 10. If they value the flower, their score would be whatever they value it at. But if they don't value the flower, maybe they want the frames. Hmm.} \\
\\ 
\texttt{Alternatively, if they value the frames more than the} \\
            }
        }
    }
    & \\ \\

    \theutterance \stepcounter{utterance}  
    & & & \multicolumn{4}{p{0.6\linewidth}}{
        \cellcolor[rgb]{0.9,0.9,0.9}{
            \makecell[{{p{\linewidth}}}]{
                \texttt{\tiny{[GM$\langle$P2]}}
                \texttt{<think>} \\
\texttt{Okay, let's see. I need to figure out how to divide the items between me and the other player. The items are 1 pen, 1 flower, and 3 frames. My values are pen:4, flower:0, frame:2. The other player's values are secret, so I don't know what they want. My goal is to maximize the total score (mine plus theirs) by making sure our proposals are complementary.} \\
\\ 
\texttt{First, I should probably ask the other player what they value most. Since I can send up to 5 messages, I can use that to get some information. Let me start by asking them which items they value the most. For example, "What do you value most: pen, flower, or frames?" That way, they can respond and I can adjust my proposal accordingly.} \\
\\ 
\texttt{If they say they value the pen, then I might want to give them the pen and take the frames and flower. But wait, I value the pen at 4, which is pretty high. If they also value the pen, maybe we need to negotiate. But if they don't value the pen, maybe I should take it. Alternatively, if they value the flower, which I don't care about, I can take the pen and frames, and let them have the flower.} \\
\\ 
\texttt{Frames are 3 in total. If both of us value frames, we need to split them. Since I get 2 per frame, 3 frames would give me 6. If the other player also values frames, maybe we can split them 2 and 1. But without knowing their value, it's hard to say. Maybe I should propose taking 2 frames and the pen, leaving 1 frame and the flower for them. But if they value the flower, that might work. If they don't, maybe they want the frames.} \\
\\ 
\texttt{Wait, the flower is only 1, and I value it at 0. So if the other player values the flower, they might want it. So maybe I should take the pen and all 3 frames, and let them have the flower. That way, my score would be 4 (pen) + 3*2 (frames) = 10. If they value the flower, their score would be whatever they value it at. But if they don't value the flower, maybe they want the frames. Hmm.} \\
\\ 
\texttt{Alternatively, if they value the frames more than the} \\
            }
        }
    }
     \\ \\

    \theutterance \stepcounter{utterance}  
    & & & \multicolumn{2}{p{0.3\linewidth}}{
        \cellcolor[rgb]{0.9,0.9,0.9}{
            \makecell[{{p{\linewidth}}}]{
                \texttt{\tiny{[GM$|$GM]}}
                \texttt{continue} \\
            }
        }
    }
    & & \\ \\

    \theutterance \stepcounter{utterance}  
    & & & \multicolumn{2}{p{0.3\linewidth}}{
        \cellcolor[rgb]{0.9,0.9,0.9}{
            \makecell[{{p{\linewidth}}}]{
                \texttt{\tiny{[GM$|$GM]}}
                \texttt{<think>} \\
\texttt{Okay, let's see. I need to figure out how to divide the items between me and the other player. The items are 1 pen, 1 flower, and 3 frames. My values are pen:4, flower:0, frame:2. The other player's values are secret, so I don't know what they want. My goal is to maximize the total score (mine plus theirs) by making sure our proposals are complementary.} \\
\\ 
\texttt{First, I should probably ask the other player what they value most. Since I can send up to 5 messages, I can use that to get some information. Let me start by asking them which items they value the most. For example, "What do you value most: pen, flower, or frames?" That way, they can respond and I can adjust my proposal accordingly.} \\
\\ 
\texttt{If they say they value the pen, then I might want to give them the pen and take the frames and flower. But wait, I value the pen at 4, which is pretty high. If they also value the pen, maybe we need to negotiate. But if they don't value the pen, maybe I should take it. Alternatively, if they value the flower, which I don't care about, I can take the pen and frames, and let them have the flower.} \\
\\ 
\texttt{Frames are 3 in total. If both of us value frames, we need to split them. Since I get 2 per frame, 3 frames would give me 6. If the other player also values frames, maybe we can split them 2 and 1. But without knowing their value, it's hard to say. Maybe I should propose taking 2 frames and the pen, leaving 1 frame and the flower for them. But if they value the flower, that might work. If they don't, maybe they want the frames.} \\
\\ 
\texttt{Wait, the flower is only 1, and I value it at 0. So if the other player values the flower, they might want it. So maybe I should take the pen and all 3 frames, and let them have the flower. That way, my score would be 4 (pen) + 3*2 (frames) = 10. If they value the flower, their score would be whatever they value it at. But if they don't value the flower, maybe they want the frames. Hmm.} \\
\\ 
\texttt{Alternatively, if they value the frames more than the} \\
            }
        }
    }
    & & \\ \\

    \theutterance \stepcounter{utterance}  
    & & \multicolumn{4}{p{0.6\linewidth}}{
        \cellcolor[rgb]{0.9,0.9,0.9}{
            \makecell[{{p{\linewidth}}}]{
                \texttt{\tiny{[P1$\langle$GM]}}
                \texttt{<think>} \\
\texttt{Okay, let's see. I need to figure out how to divide the items between me and the other player. The items are 1 pen, 1 flower, and 3 frames. My values are pen:4, flower:0, frame:2. The other player's values are secret, so I don't know what they want. My goal is to maximize the total score (mine plus theirs) by making sure our proposals are complementary.} \\
\\ 
\texttt{First, I should probably ask the other player what they value most. Since I can send up to 5 messages, I can use that to get some information. Let me start by asking them which items they value the most. For example, "What do you value most: pen, flower, or frames?" That way, they can respond and I can adjust my proposal accordingly.} \\
\\ 
\texttt{If they say they value the pen, then I might want to give them the pen and take the frames and flower. But wait, I value the pen at 4, which is pretty high. If they also value the pen, maybe we need to negotiate. But if they don't value the pen, maybe I should take it. Alternatively, if they value the flower, which I don't care about, I can take the pen and frames, and let them have the flower.} \\
\\ 
\texttt{Frames are 3 in total. If both of us value frames, we need to split them. Since I get 2 per frame, 3 frames would give me 6. If the other player also values frames, maybe we can split them 2 and 1. But without knowing their value, it's hard to say. Maybe I should propose taking 2 frames and the pen, leaving 1 frame and the flower for them. But if they value the flower, that might work. If they don't, maybe they want the frames.} \\
\\ 
\texttt{Wait, the flower is only 1, and I value it at 0. So if the other player values the flower, they might want it. So maybe I should take the pen and all 3 frames, and let them have the flower. That way, my score would be 4 (pen) + 3*2 (frames) = 10. If they value the flower, their score would be whatever they value it at. But if they don't value the flower, maybe they want the frames. Hmm.} \\
\\ 
\texttt{Alternatively, if they value the frames more than the} \\
            }
        }
    }
    & \\ \\

    \theutterance \stepcounter{utterance}  
    & \multicolumn{4}{p{0.6\linewidth}}{
        \cellcolor[rgb]{0.9,0.9,0.9}{
            \makecell[{{p{\linewidth}}}]{
                \texttt{\tiny{[P1$\rangle$GM]}}
                \texttt{<think>} \\
\texttt{Okay, let's see. I need to figure out how to divide the items between me and the other player. The items are 1 pen, 1 flower, and 3 frames. My values are pen:4, flower:0, frame:2. The other player's values are secret, so I don't know what they want. My goal is to maximize the total score (mine plus theirs) by making sure our proposals are complementary.} \\
\\ 
\texttt{First, I should probably ask the other player what they value most. Since I can send up to 5 messages, I can use that to get some information. Let me start by asking them which items they value the most. For example, "What do you value most: pen, flower, or frames?" That way, they can respond and I can adjust my proposal accordingly.} \\
\\ 
\texttt{If they say they value the pen, then I might want to give them the pen and take the frames and flower. But wait, I value the pen at 4, which is pretty high. If they also value the pen, maybe we need to negotiate. But if they don't value the pen, maybe I should take it. Alternatively, if they value the flower, which I don't care about, I can take the pen and frames, and let them have the flower.} \\
\\ 
\texttt{Frames are 3 in total. If both of us value frames, we need to split them. Since I get 2 per frame, 3 frames would give me 6. If the other player also values frames, maybe we can split them 2 and 1. But without knowing their value, it's hard to say. Maybe I should propose taking 2 frames and the pen, leaving 1 frame and the flower for them. But if they value the flower, that might work. If they don't, maybe they want the frames.} \\
\\ 
\texttt{Wait, the flower is only 1, and I value it at 0. So if the other player values the flower, they might want it. So maybe I should take the pen and all 3 frames, and let them have the flower. That way, my score would be 4 (pen) + 3*2 (frames) = 10. If they value the flower, their score would be whatever they value it at. But if they don't value the flower, maybe they want the frames. Hmm.} \\
\\ 
\texttt{Alternatively, if they value the frames more than the} \\
            }
        }
    }
    & & \\ \\

    \theutterance \stepcounter{utterance}  
    & & & \multicolumn{2}{p{0.3\linewidth}}{
        \cellcolor[rgb]{0.9,0.9,0.9}{
            \makecell[{{p{\linewidth}}}]{
                \texttt{\tiny{[GM$|$GM]}}
                \texttt{continue} \\
            }
        }
    }
    & & \\ \\

    \theutterance \stepcounter{utterance}  
    & & & \multicolumn{2}{p{0.3\linewidth}}{
        \cellcolor[rgb]{0.9,0.9,0.9}{
            \makecell[{{p{\linewidth}}}]{
                \texttt{\tiny{[GM$|$GM]}}
                \texttt{<think>} \\
\texttt{Okay, let's see. I need to figure out how to divide the items between me and the other player. The items are 1 pen, 1 flower, and 3 frames. My values are pen:4, flower:0, frame:2. The other player's values are secret, so I don't know what they want. My goal is to maximize the total score (mine plus theirs) by making sure our proposals are complementary.} \\
\\ 
\texttt{First, I should probably ask the other player what they value most. Since I can send up to 5 messages, I can use that to get some information. Let me start by asking them which items they value the most. For example, "What do you value most: pen, flower, or frames?" That way, they can respond and I can adjust my proposal accordingly.} \\
\\ 
\texttt{If they say they value the pen, then I might want to give them the pen and take the frames and flower. But wait, I value the pen at 4, which is pretty high. If they also value the pen, maybe we need to negotiate. But if they don't value the pen, maybe I should take it. Alternatively, if they value the flower, which I don't care about, I can take the pen and frames, and let them have the flower.} \\
\\ 
\texttt{Frames are 3 in total. If both of us value frames, we need to split them. Since I get 2 per frame, 3 frames would give me 6. If the other player also values frames, maybe we can split them 2 and 1. But without knowing their value, it's hard to say. Maybe I should propose taking 2 frames and the pen, leaving 1 frame and the flower for them. But if they value the flower, that might work. If they don't, maybe they want the frames.} \\
\\ 
\texttt{Wait, the flower is only 1, and I value it at 0. So if the other player values the flower, they might want it. So maybe I should take the pen and all 3 frames, and let them have the flower. That way, my score would be 4 (pen) + 3*2 (frames) = 10. If they value the flower, their score would be whatever they value it at. But if they don't value the flower, maybe they want the frames. Hmm.} \\
\\ 
\texttt{Alternatively, if they value the frames more than the} \\
            }
        }
    }
    & & \\ \\

    \theutterance \stepcounter{utterance}  
    & & \multicolumn{4}{p{0.6\linewidth}}{
        \cellcolor[rgb]{0.9,0.9,0.9}{
            \makecell[{{p{\linewidth}}}]{
                \texttt{\tiny{[GM$\rangle$P2]}}
                \texttt{<think>} \\
\texttt{Okay, let's see. I need to figure out how to divide the items between me and the other player. The items are 1 pen, 1 flower, and 3 frames. My values are pen:4, flower:0, frame:2. The other player's values are secret, so I don't know what they want. My goal is to maximize the total score (mine plus theirs) by making sure our proposals are complementary.} \\
\\ 
\texttt{First, I should probably ask the other player what they value most. Since I can send up to 5 messages, I can use that to get some information. Let me start by asking them which items they value the most. For example, "What do you value most: pen, flower, or frames?" That way, they can respond and I can adjust my proposal accordingly.} \\
\\ 
\texttt{If they say they value the pen, then I might want to give them the pen and take the frames and flower. But wait, I value the pen at 4, which is pretty high. If they also value the pen, maybe we need to negotiate. But if they don't value the pen, maybe I should take it. Alternatively, if they value the flower, which I don't care about, I can take the pen and frames, and let them have the flower.} \\
\\ 
\texttt{Frames are 3 in total. If both of us value frames, we need to split them. Since I get 2 per frame, 3 frames would give me 6. If the other player also values frames, maybe we can split them 2 and 1. But without knowing their value, it's hard to say. Maybe I should propose taking 2 frames and the pen, leaving 1 frame and the flower for them. But if they value the flower, that might work. If they don't, maybe they want the frames.} \\
\\ 
\texttt{Wait, the flower is only 1, and I value it at 0. So if the other player values the flower, they might want it. So maybe I should take the pen and all 3 frames, and let them have the flower. That way, my score would be 4 (pen) + 3*2 (frames) = 10. If they value the flower, their score would be whatever they value it at. But if they don't value the flower, maybe they want the frames. Hmm.} \\
\\ 
\texttt{Alternatively, if they value the frames more than the} \\
            }
        }
    }
    & \\ \\

    \theutterance \stepcounter{utterance}  
    & & & \multicolumn{4}{p{0.6\linewidth}}{
        \cellcolor[rgb]{0.9,0.9,0.9}{
            \makecell[{{p{\linewidth}}}]{
                \texttt{\tiny{[GM$\langle$P2]}}
                \texttt{<think>} \\
\texttt{Okay, let's see. I need to figure out how to divide the items between me and the other player. The items are 1 pen, 1 flower, and 3 frames. My values are pen:4, flower:0, frame:2. The other player's values are secret, so I don't know what they want. My goal is to maximize the total score (mine plus theirs) by making sure our proposals are complementary.} \\
\\ 
\texttt{First, I should probably ask the other player what they value most. Since I can send up to 5 messages, I can use that to get some information. Let me start by asking them which items they value the most. For example, "What do you value most: pen, flower, or frames?" That way, they can respond and I can adjust my proposal accordingly.} \\
\\ 
\texttt{If they say they value the pen, then I might want to give them the pen and take the frames and flower. But wait, I value the pen at 4, which is pretty high. If they also value the pen, maybe we need to negotiate. But if they don't value the pen, maybe I should take it. Alternatively, if they value the flower, which I don't care about, I can take the pen and frames, and let them have the flower.} \\
\\ 
\texttt{Frames are 3 in total. If both of us value frames, we need to split them. Since I get 2 per frame, 3 frames would give me 6. If the other player also values frames, maybe we can split them 2 and 1. But without knowing their value, it's hard to say. Maybe I should propose taking 2 frames and the pen, leaving 1 frame and the flower for them. But if they value the flower, that might work. If they don't, maybe they want the frames.} \\
\\ 
\texttt{Wait, the flower is only 1, and I value it at 0. So if the other player values the flower, they might want it. So maybe I should take the pen and all 3 frames, and let them have the flower. That way, my score would be 4 (pen) + 3*2 (frames) = 10. If they value the flower, their score would be whatever they value it at. But if they don't value the flower, maybe they want the frames. Hmm.} \\
\\ 
\texttt{Alternatively, if they value the frames more than the} \\
            }
        }
    }
     \\ \\

    \theutterance \stepcounter{utterance}  
    & & & \multicolumn{2}{p{0.3\linewidth}}{
        \cellcolor[rgb]{0.9,0.9,0.9}{
            \makecell[{{p{\linewidth}}}]{
                \texttt{\tiny{[GM$|$GM]}}
                \texttt{continue} \\
            }
        }
    }
    & & \\ \\

    \theutterance \stepcounter{utterance}  
    & & & \multicolumn{2}{p{0.3\linewidth}}{
        \cellcolor[rgb]{0.9,0.9,0.9}{
            \makecell[{{p{\linewidth}}}]{
                \texttt{\tiny{[GM$|$GM]}}
                \texttt{<think>} \\
\texttt{Okay, let's see. I need to figure out how to divide the items between me and the other player. The items are 1 pen, 1 flower, and 3 frames. My values are pen:4, flower:0, frame:2. The other player's values are secret, so I don't know what they want. My goal is to maximize the total score (mine plus theirs) by making sure our proposals are complementary.} \\
\\ 
\texttt{First, I should probably ask the other player what they value most. Since I can send up to 5 messages, I can use that to get some information. Let me start by asking them which items they value the most. For example, "What do you value most: pen, flower, or frames?" That way, they can respond and I can adjust my proposal accordingly.} \\
\\ 
\texttt{If they say they value the pen, then I might want to give them the pen and take the frames and flower. But wait, I value the pen at 4, which is pretty high. If they also value the pen, maybe we need to negotiate. But if they don't value the pen, maybe I should take it. Alternatively, if they value the flower, which I don't care about, I can take the pen and frames, and let them have the flower.} \\
\\ 
\texttt{Frames are 3 in total. If both of us value frames, we need to split them. Since I get 2 per frame, 3 frames would give me 6. If the other player also values frames, maybe we can split them 2 and 1. But without knowing their value, it's hard to say. Maybe I should propose taking 2 frames and the pen, leaving 1 frame and the flower for them. But if they value the flower, that might work. If they don't, maybe they want the frames.} \\
\\ 
\texttt{Wait, the flower is only 1, and I value it at 0. So if the other player values the flower, they might want it. So maybe I should take the pen and all 3 frames, and let them have the flower. That way, my score would be 4 (pen) + 3*2 (frames) = 10. If they value the flower, their score would be whatever they value it at. But if they don't value the flower, maybe they want the frames. Hmm.} \\
\\ 
\texttt{Alternatively, if they value the frames more than the} \\
            }
        }
    }
    & & \\ \\

    \theutterance \stepcounter{utterance}  
    & & \multicolumn{4}{p{0.6\linewidth}}{
        \cellcolor[rgb]{0.9,0.9,0.9}{
            \makecell[{{p{\linewidth}}}]{
                \texttt{\tiny{[P1$\langle$GM]}}
                \texttt{<think>} \\
\texttt{Okay, let's see. I need to figure out how to divide the items between me and the other player. The items are 1 pen, 1 flower, and 3 frames. My values are pen:4, flower:0, frame:2. The other player's values are secret, so I don't know what they want. My goal is to maximize the total score (mine plus theirs) by making sure our proposals are complementary.} \\
\\ 
\texttt{First, I should probably ask the other player what they value most. Since I can send up to 5 messages, I can use that to get some information. Let me start by asking them which items they value the most. For example, "What do you value most: pen, flower, or frames?" That way, they can respond and I can adjust my proposal accordingly.} \\
\\ 
\texttt{If they say they value the pen, then I might want to give them the pen and take the frames and flower. But wait, I value the pen at 4, which is pretty high. If they also value the pen, maybe we need to negotiate. But if they don't value the pen, maybe I should take it. Alternatively, if they value the flower, which I don't care about, I can take the pen and frames, and let them have the flower.} \\
\\ 
\texttt{Frames are 3 in total. If both of us value frames, we need to split them. Since I get 2 per frame, 3 frames would give me 6. If the other player also values frames, maybe we can split them 2 and 1. But without knowing their value, it's hard to say. Maybe I should propose taking 2 frames and the pen, leaving 1 frame and the flower for them. But if they value the flower, that might work. If they don't, maybe they want the frames.} \\
\\ 
\texttt{Wait, the flower is only 1, and I value it at 0. So if the other player values the flower, they might want it. So maybe I should take the pen and all 3 frames, and let them have the flower. That way, my score would be 4 (pen) + 3*2 (frames) = 10. If they value the flower, their score would be whatever they value it at. But if they don't value the flower, maybe they want the frames. Hmm.} \\
\\ 
\texttt{Alternatively, if they value the frames more than the} \\
\\ 
\\ 
\texttt{The time is up. It is now your turn to submit a secret proposal.} \\
            }
        }
    }
    & \\ \\

    \theutterance \stepcounter{utterance}  
    & \multicolumn{4}{p{0.6\linewidth}}{
        \cellcolor[rgb]{0.9,0.9,0.9}{
            \makecell[{{p{\linewidth}}}]{
                \texttt{\tiny{[P1$\rangle$GM]}}
                \texttt{<think>} \\
\texttt{Okay, let's see. I need to make a secret proposal here. The items are 1 pen, 1 flower, and 3 frames. My values are pen:4, flower:0, frame:2. The other player's values are unknown. My goal is to maximize the total score (mine plus theirs) by making sure our proposals are complementary.} \\
\\ 
\texttt{Since I don't know the other player's values, I have to make an educated guess. The flower is only 1, and I value it at 0, so maybe the other player values it. If they do, giving them the flower would be good. The pen is valuable to me (4 points), so I probably want to take it. Frames are 3, each worth 2 to me, so 6 total. If the other player also values frames, we need to split them. But without knowing their values, I should assume they might want some frames. Let's say I take the pen and 2 frames (total 4 + 4 = 8), and leave 1 frame and the flower for them. That way, if they value the flower or frames, they can take those. My proposal would be {[}Proposal: 1 pen, 2 frame{]}. That leaves 1 frame and 1 flower for them. If their proposal is for those, it's complementary. If not, maybe they want the pen, but I already took it. Hmm, but if they also want the pen, our proposals would conflict. But since I value it higher, maybe it's better for me to take it. Let's go with that.} \\
\texttt{</think>} \\
\\ 
\texttt{{[}Proposal: 1 pen, 2 frame{]}} \\
            }
        }
    }
    & & \\ \\

    \theutterance \stepcounter{utterance}  
    & & & \multicolumn{2}{p{0.3\linewidth}}{
        \cellcolor[rgb]{0.9,0.9,0.9}{
            \makecell[{{p{\linewidth}}}]{
                \texttt{\tiny{[GM$|$GM]}}
                \texttt{proposal} \\
            }
        }
    }
    & & \\ \\

    \theutterance \stepcounter{utterance}  
    & & & \multicolumn{2}{p{0.3\linewidth}}{
        \cellcolor[rgb]{0.9,0.9,0.9}{
            \makecell[{{p{\linewidth}}}]{
                \texttt{\tiny{[GM$|$GM]}}
                [1, 0, 2]
            }
        }
    }
    & & \\ \\

    \theutterance \stepcounter{utterance}  
    & & \multicolumn{4}{p{0.6\linewidth}}{
        \cellcolor[rgb]{0.9,0.9,0.9}{
            \makecell[{{p{\linewidth}}}]{
                \texttt{\tiny{[GM$\rangle$P2]}}
                \texttt{The other player has submitted a secret proposal. It is now your turn to submit a proposal of your own.} \\
            }
        }
    }
    & \\ \\

    \theutterance \stepcounter{utterance}  
    & & & \multicolumn{4}{p{0.6\linewidth}}{
        \cellcolor[rgb]{0.9,0.9,0.9}{
            \makecell[{{p{\linewidth}}}]{
                \texttt{\tiny{[GM$\langle$P2]}}
                \texttt{<think>} \\
\texttt{Okay, the other player has already submitted a secret proposal, and now I need to submit mine. Since I don't know their proposal, I have to make sure mine is complementary. Let me think about the possible scenarios.} \\
\\ 
\texttt{First, the items are 1 pen, 1 flower, and 3 frames. My values are pen:4, flower:0, frame:2. The other player's values are unknown. My goal is to maximize the total score (mine plus theirs) by ensuring our proposals don't overlap.} \\
\\ 
\texttt{If the other player asked for the pen, I should take the frames and flower. But since I don't value the flower, maybe they want it. If they took the pen, I can take all 3 frames (giving me 6 points) and the flower (0 for me). Total for me would be 6, and if they value the pen, their score would be whatever they have. But if they also want frames, maybe they took some, so I need to adjust.} \\
\\ 
\texttt{Alternatively, if they proposed taking the flower, I can take the pen and frames. That would give me 4 + 6 = 10 points. If they value the flower, their score would be their value for it, and mine is 10. That seems good.} \\
\\ 
\texttt{If they took some frames, say 2 frames, then I can take the remaining 1 frame, pen, and flower. My score would be 4 + 2 + 0 = 6. But maybe they took 1 frame, so I take 2 frames, pen, and flower. That's 4 + 4 + 0 = 8. Hmm.} \\
\\ 
\texttt{Wait, the key is to make sure our proposals don't overlap. Since I don't know what they took, I need to propose a division that covers all items without conflict. The safest way is to take the items I value most and leave the rest. Since I value the pen at 4 and frames at 2 each, I should prioritize taking the pen and as many frames as possible. The flower is worthless to me, so I can leave it for them.} \\
\\ 
\texttt{So my proposal should be: 1 pen and 3 frames. That way, I get 4 + 3*2 = 10 points. If the other player's proposal includes the flower, then our proposals are complementary. If they also wanted the pen or frames, there might} \\
            }
        }
    }
     \\ \\

    \theutterance \stepcounter{utterance}  
    & & & \multicolumn{2}{p{0.3\linewidth}}{
        \cellcolor[rgb]{0.9,0.9,0.9}{
            \makecell[{{p{\linewidth}}}]{
                \texttt{\tiny{[GM$|$GM]}}
                \texttt{continue} \\
            }
        }
    }
    & & \\ \\

    \theutterance \stepcounter{utterance}  
    & & & \multicolumn{2}{p{0.3\linewidth}}{
        \cellcolor[rgb]{0.9,0.9,0.9}{
            \makecell[{{p{\linewidth}}}]{
                \texttt{\tiny{[GM$|$GM]}}
                \texttt{player was instructed to make a proposal} \\
            }
        }
    }
    & & \\ \\

\end{supertabular}
}

\end{document}
